\begin{helper}
Fix a two-bite sample \(\omega\).  Suppose every lifted neighborhood has
size at most \(2pn\), the red-red and blue-blue lifted intersections are each
at most \(100(\log n)^3\), and the collision event supplied by
\(\noderef{OppositeProjectionCollisionBound}\) holds for this sample.  Then
the two opposite-projection clauses of \(\noderef{FiberAndDegreeEvent}\) hold:
for all distinct red vertices \(r,s\),
\[
|\pi_B(N_r)\cap\pi_B(N_s)|\le 2\cdot 100(\log n)^3,
\]
and for all distinct blue vertices \(b,c\),
\[
|\pi_R(N_b)\cap\pi_R(N_c)|\le 2\cdot 100(\log n)^3.
\]
\end{helper}
\begin{proof}
The collision event says that the size bound implies both same-color
opposite-projection collision estimates.  Applying it to the assumed lifted
size bound gives, for distinct red vertices \(r,s\),
\[
|\pi_B(N_r)\cap\pi_B(N_s)|
\le |N_r\cap N_s|+100(\log n)^3.
\]
The assumed red-red lifted-intersection bound gives
\(|N_r\cap N_s|\le100(\log n)^3\), so the displayed inequality is at most
\(2\cdot100(\log n)^3\).  The blue statement is identical, using the blue-blue
lifted-intersection bound and the blue half of the collision event.
\end{proof}
